% ============================================================================
\chapter{Le navigateur : HTML, CSS et le DOM}
\label{ch:js-dom}
% ============================================================================

\begin{objectif}
\begin{itemize}
  \item Poser une page HTML correcte, accessible, et y brancher du JavaScript proprement.
  \item S\'electionner, cr\'eer, modifier, supprimer des noeuds sans casser les performances.
  \item Comprendre le rendu (reflow/repaint) et \'ecrire du code qui reste fluide \`a 60 fps.
  \item Utiliser les API modernes~: \texttt{dataset}, \texttt{classList}, \texttt{dialog},
        \texttt{details}, \texttt{IntersectionObserver}.
  \item Livrer un premier mini-projet complet, sans framework.
\end{itemize}
\end{objectif}

\section{Le squelette d'une page moderne}

\begin{codeplain}
<!DOCTYPE html>
<html lang="fr">
<head>
  <meta charset="utf-8">
  <meta name="viewport" content="width=device-width, initial-scale=1">
  <title>Liste de courses -- Nour</title>
  <meta name="description" content="Une liste de courses qui persiste dans le navigateur.">
  <meta name="color-scheme" content="light dark">
  <link rel="icon" href="/favicon.svg" type="image/svg+xml">
  <link rel="stylesheet" href="style.css">
  <script src="app.js" defer></script>            <!-- defer = execute apres le parse du HTML -->
</head>
<body>
  <a class="sr-only" href="#contenu">Aller au contenu</a>
  <header>
    <h1>Mes courses</h1>
  </header>
  <main id="contenu">
    <form id="form-ajout" autocomplete="off">
      <label for="nom">Article</label>
      <input id="nom" name="nom" required maxlength="60" placeholder="Riz, tomates...">
      <label for="qte">Quantit\'e</label>
      <input id="qte" name="qte" type="number" min="1" max="99" value="1">
      <button type="submit">Ajouter</button>
    </form>
    <ul id="liste" aria-live="polite"></ul>
    <template id="tpl-article">
      <li class="article">
        <input type="checkbox" class="check">
        <span class="nom"></span>
        <button class="suppr" type="button" aria-label="Supprimer">x</button>
      </li>
    </template>
  </main>
  <footer><small><span id="compteur"></span></small></footer>
</body>
</html>
\end{codeplain}

Quatre d\'etails qui font la diff\'erence~: \cmd{lang="fr"} (prononciation, c\'esure, recherche),
la balise \verb|viewport| (sinon ton site reste \og ~zoom\'e~ \fg{} sur mobile), \verb|defer|
(le script attend le HTML, sans bloquer le parsing) et le \verb|<template>| (du HTML stock\'e,
inerte, clonable \`a la demande -- l'anti-\verb|innerHTML| par excellence).

\begin{tabz}{lX}
  \headrow \thd{Attribut de script} \& \thd{Comportement} \\
  \midrule
  \verb|<script src="a.js">| \& bloque le parsing jusqu'au t\'el\'echargement et \`a l'ex\'ecution. \\
  \verb|defer| \& t\'el\'echargement en parall\`ele, ex\'ecution dans l'ordre, apr\`es le parsing. \textbf{Le choix par d\'efaut pour un fichier externe.} \\
  \verb|async| \& t\'el\'echargement en parall\`ele, ex\'ecution d\`es que pr\^et, ordre al\'eatoire. Pour la pub, l'analytics. \\
  \verb|type="module"| \& diff\'er\'e par d\'efaut, port\'ee de module, \verb|import| disponibles. \textbf{Le choix pour du JS en modules.} \\
  \verb|inlin\'e <script>| \& ex\'ecut\'e imm\'ediatement~: \`a placer en \emph{fin de body} si n\'ecessaire. \\
\end{tabz}

\section{S\'electionner}

\begin{codejs}
const box = document.querySelector(".boite");       // 1er element correspondant
const tous = document.querySelectorAll(".article");  // NodeList statique (tableau-like)
const parId = document.getElementById("liste");      // le plus rapide, a privilegier si tu as un id
const parClasse = document.getElementsByClassName("article");   // HTMLCollection VIVANTE (piege)
const parNom = document.getElementsByTagName("li");

// Selections structurelles
const parent = box.closest("main");          // remonte jusqu'au premier ancetre qui matche
const suivant = box.nextElementSibling;
const enfants = [...box.children];           // HTMLCollection -> vrai tableau
const dansForm = document.querySelectorAll("form input:not([type=hidden])");
const visible = [...document.querySelectorAll(".article")].filter((el) => el.checkVisibility());
box.matches("main > .boite");                 // test de selecteur

// Les selecteurs utiles a connaitre par coeur
// :not(), :is(h1,h2,h3), :where(), :has(> input:checked), :nth-child(2n+1), :scope, ::selection
document.querySelectorAll("li:has(.check:checked)");     // les lignes cochees (CSS + JS, tres puissant)
\end{codejs}

\begin{note}
\texttt{querySelectorAll} renvoie une copie fig\'ee~: si tu ajoutes un \'el\'ement apr\`es, il n'y sera
pas. \texttt{getElementsByClassName} renvoie une collection \emph{vivante}~: la modifier pendant
une boucle produit des sauts d'index. Dans le doute~: \verb|[...document.querySelectorAll(...)]|.
\end{note}

\section{Lire et modifier un \'el\'ement}

\begin{codejs}
const li = document.querySelector(".article");

// TEXTE : textContent securise (les <b> sont echappes), innerHTML non (XSS)
li.textContent = "Riz <5kg>";              // affiche le texte literal, y compris les chevrons
li.innerHTML = "<b>Riz</b> 5kg";           // parse du HTML : UNIQUEMENT avec du contenu de confiance
li.insertAdjacentHTML("beforeend", "<i>bio</i>");  // le plus rapide des trois pour inserer du HTML

// CLASSES ET ATTRIBUTS
li.classList.add("actif"); li.classList.remove("cache");
li.classList.toggle("fait", coche);         // 2e argument = booleen : la forme propre
li.classList.replace("petit", "grand");
li.classList.contains("fait");
li.className = "a b c";                     // ecrase tout : a eviter
li.setAttribute("aria-label", "Article riz");
li.getAttribute("data-id");                 // "42" (toujours une chaine)
li.id = "article-42";
li.hidden = true;                           // equivalent de [hidden]

// DONNEES UTILISATEUR : dataset (kebab-case <-> camelCase)
// <li data-id="42" data-categorie="feculent">
li.dataset.id;                // "42"
li.dataset.categorie = "feculent";
li.dataset.etatCoche = "1";   // data-etat-coche

// STYLE : prefere une classe CSS a une propriete inline
li.style.backgroundColor = "#eef";          // camelCase en JS, kebab-case en CSS
li.style.setProperty("--hauteur", "40px");  // variable CSS
getComputedStyle(li).fontSize;              // "16px" (la valeur calculee, pas celle ecrite)
li.style.cssText = "color:red;padding:4px"; // remplace tout le style inline

// DIMENSIONS (attention : forcent un reflow -- a lire par lots)
li.offsetWidth; li.offsetHeight; li.getBoundingClientRect();
document.documentElement.scrollTop;
\end{codejs}

\subsection{Cr\'eer, ins\'erer, supprimer, d\'eplacer}

\begin{codejs}
// La voie royale : template + clone (rapide, sans parse, sans injection)
const tpl = document.getElementById("tpl-article");
const creerArticle = ({ nom, qte }) => {
  const frag = tpl.content.cloneNode(true);            // true = avec les descendants
  frag.querySelector(".nom").textContent = `${nom} x${qte}`;
  frag.querySelector(".suppr").dataset.id = nom;
  return frag;
};
document.getElementById("liste").append(creerArticle({ nom: "Riz", qte: 2 }));

// Creation manuelle
const el = document.createElement("li");
el.className = "article";
el.append("texte ", document.createElement("b"), " et noeud");   // append accepte textes et noeuds
el.prepend(puce);
el.replaceChildren(...listeFils);        // remplace TOUS les enfants en une fois : tres efficace
el.before(sibling); el.after(sibling); el.replaceWith(autre);
el.remove();                              // le fils se detache du parent ; l'objet reste utilisable

// Ou par HTML (rappel : seulement pour du contenu de confiance)
const conteneur = document.getElementById("liste");
conteneur.insertAdjacentHTML("beforeend", htmlSecurise);
conteneur.children.length;
\end{codejs}

\etape{\'Ecrire une fonction de nettoyage si tu dois ins\'erer du HTML}
\begin{codejs}
const echapper = (s) => String(s)
  .replaceAll("&", "&amp;").replaceAll("<", "&lt;").replaceAll(">", "&gt;")
  .replaceAll('"', "&quot;").replaceAll("'", "&#39;");
const rendu = (t) => `<strong>${echapper(t)}</strong>`;   // le texte est sur, la structure est la tienne
\end{codejs}
En production, ne r\'einvente pas ce filtre~: \link{DOMPurify}{https://github.com/cure53/DOMPurify}
(\cmd{npm i dompurify}) g\`ere les cas tordus \verb|<svg>|, les gestionnaires inline, les URL
\verb|javascript:|.

\section{Le mod\`ele de l'arbre et les noeuds}

\begin{codejs}
document            // Document
 .documentElement    // <html>
 .firstElementChild  // <head>
document.body; document.head; document.title;
document.activeElement;                       // l'element qui a le focus
document.getSelection().toString();           // texte selectionne
node.parentNode; node.parentElement;         // parent, en incluant les noeuds de texte
node.childNodes; node.children;               // idem : tous vs elements uniquement
node.nextSibling / node.nextElementSibling;
node.isConnected;                              // encore dans le document ?
node.contains(other);
node.innerText;                                // comme textContent mais respectant le CSS (lent !)
\end{codejs}

\begin{piege}
\verb|innerText| d\'eclenche un calcul de style complet sur l'\'el\'ement et ses descendants.
Sur une liste de mille lignes, c'est plusieurs secondes de g\^ach\'ees. Utilise
\verb|textContent|~; garde \verb|innerText| pour reproduire l'affichage visible (sauts de ligne,
\verb|display:none|).
\end{piege}

\section{Performance : ce que le navigateur paie}

Le pipeline est~: JavaScript $\rightarrow$ styles (recalcul) $\rightarrow$ layout (g\'eom\'etrie)
$\rightarrow$ paint (pixels) $\rightarrow$ composite (GPU). Tout changement de
\texttt{width/top/left/height} relance layout~; tout changement de
\texttt{transform/opacity} peut rester sur le compositeur~: c'est le chemin rapide.

\begin{codejs}
// ANTI-PATTERN : lecture puis ecriture = reflow a chaque tour (layout thrashing)
for (const el of elements) {
  el.style.height = el.offsetHeight + 10 + "px";      // lit la geometrie (reflow) puis ecrit
}

// CORRECT : lire tout, puis ecrire tout
const hauteurs = elements.map((el) => el.offsetHeight);
elements.forEach((el, i) => { el.style.height = `${hauteurs[i] + 10}px`; });

// ENCORE MIEUX : laisser le CSS faire le travail
el.classList.add("agrandi");
// Animer uniquement transform et opacity
el.animate(
  [{ opacity: 0, transform: "translateY(8px)" }, { opacity: 1, transform: "none" }],
  { duration: 200, easing: "cubic-bezier(.2,.8,.2,1)", fill: "both" }
);

// Grouper les mutations importantes dans une frame
requestAnimationFrame(() => { miseAJourLourde(); });

// Observer les changements de DOM sans polluer : MutationObserver
const obs = new MutationObserver((mutations) => {
  for (const m of mutations) {
    if (m.type === "childList") console.log("ajoutes :", [...m.addedNodes].length, "retires :", m.removedNodes.length);
    if (m.type === "attributes") console.log("attribut", m.attributeName, "sur", m.target);
  }
});
obs.observe(document.getElementById("liste"), { childList: true, subtree: true, attributes: true });
obs.disconnect();

// Le lazy-loading par IntersectionObserver : la fin du scroll listener couteux
const io = new IntersectionObserver((entrees) => {
  for (const e of entrees) if (e.isIntersecting) { chargerImage(e.target); io.unobserve(e.target); }
}, { rootMargin: "200px", threshold: 0.1 });
document.querySelectorAll("img[data-src]").forEach((i) => io.observe(i));
\end{codejs}

\begin{note}
Le navigateur regroupe lui-m\^eme les \'ecritures d'un m\^eme bloc de JavaScript~: les \'ecritures
\emph{pures} ne causent pas de reflow imm\'ediat. Ce qui co\^ute, c'est \textbf{lire une valeur de
g\'eom\'etrie} juste apr\`es avoir \'ecrit. Donc~: on lit, puis on \'ecrit. Jamais en alternance.
\end{note}

\section{Composants natifs que tout le monde r\'einvente mal}

\begin{codejs}
// Boite de dialogue : gestion du focus, fermeture par Echap, piege modal inclus
const dlg = document.querySelector("dialog");
dlg.showModal();
dlg.returnValue = "ok";
dlg.addEventListener("close", () => console.log(dlg.returnValue));

// Details/summary : accord\'eon accessible sans une ligne de JS
// <details name="faq" open><summary>Question</summary><p>Reponse</p></details>
// "name" cree un groupe exclusif (un seul ouvert) : natif depuis 2023.

// Popover : menus, tooltips, sans lib
// <button popovertarget="menu">Menu</button><div id="menu" popover>...</div>

// Color/Date/Range : le controleur natif est deja accessible et mobile-friendly
// <input type="date" min="2026-01-01"> <input type="color"> <input type="range">

// Scroll fluide et ancrage
document.querySelector("#tarif").scrollIntoView({ behavior: "smooth", block: "start" });
// CSS : html { scroll-behavior: smooth }  ou  scroll-margin-top: 80px sur la cible

// Historique et URL propre (le SPA tout entier tient la-dedans)
history.pushState({ vue: "article" }, "", "/article/42");
window.addEventListener("popstate", (e) => rendre(e.state?.vue ?? "accueil"));
\end{codejs}

\section{Accessibilit\'e : trente minutes qui changent la vie de tes utilisateurs}

\begin{enumerate}
  \item Un HTML s\'emantique (\verb|main|, \verb|nav|, \verb|button|, \verb|h1..h6|) avant
        n'importe quel ARIA. \verb|<div onclick>| n'est pas un bouton~: ni focus, ni Echap,
        ni annonce au lecteur d'\'ecran.
  \item Chaque champ a un \verb|<label for>|~; les ic\^ones seules ont un \verb|aria-label|.
  \item \verb|aria-live="polite"| sur une zone qui change (compteur, r\'esultat de recherche,
        message d'erreur)~; \verb|aria-busy| pendant un chargement.
  \item Ne retire jamais l'outline de focus sans la remplacer par un style visible.
  \item Contraste minimal 4{,}5:1 pour le texte courant, 3:1 pour le gros texte
        (\link{WebAIM~Contrast~Checker}{https://webaim.org/resources/contrastchecker/}).
  \item Navigation clavier compl\`ete~: \keys{Tab} partout, \keys{Echap} ferme,
        \keys{Enter}/\keys{Space} activent, fl\`eches dans les listes.
\end{enumerate}

\begin{codejs}
// Focus piege pour une modale custom (si tu ne utilises pas <dialog>)
function piegerFocus(conteneur) {
  const selecteur = 'a[href], button:not([disabled]), input, select, textarea, [tabindex]:not([tabindex="-1"])';
  const premier = conteneur.querySelector(selecteur);
  const focusables = () => [...conteneur.querySelectorAll(selecteur)].filter((el) => el.offsetParent !== null);
  conteneur.addEventListener("keydown", (e) => {
    if (e.key !== "Tab") return;
    const els = focusables();
    const idx = els.indexOf(document.activeElement);
    if (e.shiftKey && (idx === 0 || idx === -1)) { els.at(-1).focus(); e.preventDefault(); }
    else if (!e.shiftKey && idx === els.length - 1) { els[0].focus(); e.preventDefault(); }
  });
  premier?.focus();
  return () => dernierElement?.focus();       // retablir le focus a la fermeture
}
\end{codejs}

\section{Mini-projet complet --- liste de courses persistante}

Le fichier \fichier{app.js} branch\'e sur le HTML de la \cref{ch:js-dom}~; tout est comment\'e,
et tu dois le retaper int\'egralement.

\begin{codejs}
// app.js
const stock = {
  lire: () => { try { return JSON.parse(localStorage.getItem("courses") ?? "[]"); } catch { return []; } },
  ecrire: (v) => localStorage.setItem("courses", JSON.stringify(v)),
};

const $ = (s, r = document) => r.querySelector(s);
const $$ = (s, r = document) => [...r.querySelectorAll(s)];

let articles = stock.lire();
let filtre = "tous";

function rendre() {
  const ul = $("#liste");
  const visibles = articles.filter((a) =>
    filtre === "tous" ? true : filtre === "coches" ? a.fait : !a.fait
  );
  const frag = document.createDocumentFragment();
  const tpl = $("#tpl-article");

  for (const a of visibles) {
    const n = tpl.content.cloneNode(true);
    const li = n.querySelector("li");
    li.dataset.id = a.id;
    li.classList.toggle("fait", a.fait);
    n.querySelector(".nom").textContent = a.nom;
    const check = n.querySelector(".check");
    check.checked = a.fait;
    frag.append(li.parentNode ? n : n);
  }
  ul.replaceChildren(frag);
  $("#compteur").textContent =
    `${articles.filter((a) => !a.fait).length} article(s) a acheter sur ${articles.length}`;
}

function ajouter(nom, qte) {
  nom = nom.trim();
  if (!nom) return alert("Nom requis");
  articles = [{ id: crypto.randomUUID(), nom, qte: Number(qte) || 1, fait: false }, ...articles];
  stock.ecrire(articles);
  rendre();
}

// Deleguer sur la liste : un seul ecouteur, meme pour les elements crees plus tard
$("#liste").addEventListener("click", (e) => {
  const li = e.target.closest("li[data-id]");
  if (!li) return;
  const id = li.dataset.id;
  if (e.target.matches(".check")) {
    articles = articles.map((a) => (a.id === id ? { ...a, fait: !a.fait } : a));
  } else if (e.target.matches(".suppr")) {
    articles = articles.filter((a) => a.id !== id);
  } else return;
  stock.ecrire(articles);
  rendre();
});

$("#form-ajout").addEventListener("submit", (e) => {
  e.preventDefault();
  const f = new FormData(e.target);
  ajouter(f.get("nom"), f.get("qte"));
  e.target.reset();
  $("#nom").focus();
});

$$(".filtre").forEach((b) =>
  b.addEventListener("click", () => {
    filtre = b.dataset.filtre;
    $$(".filtre").forEach((x) => x.setAttribute("aria-pressed", String(x === b)));
    rendre();
  })
);

rendre();
\end{codejs}

\begin{exercice}
\begin{enumerate}
  \item Ajoute la r\'eordonnancement par glisser-d\'eposer (\verb|draggable="true"|,
        \verb|dragstart|, \verb|dragover|, \verb|drop|) et persiste l'ordre.
  \item Ajoute un champ de recherche avec \verb|debounce|~: filtrer sur le nom, en ignorant
        accents et casse, et surligner la partie correspondante avec \verb|<mark>|.
  \item Ajoute un mode sombre~: lire \verb|prefers-color-scheme| via
        \verb|matchMedia|, stocker le choix de l'utilisateur, basculer la classe
        \verb|dark| sur \verb|<html>|.
  \item Ajoute l'annulation~: un tableau \verb|historique| des 20 derniers \'etats, et les
        raccourcis \keys{Ctrl+Z} / \keys{Ctrl+Shift+Z}.
  \item Rends le composant \emph{r\'eutilisable}~: \verb|creerListe(elise, { stock })| qui
        accepte n'importe quel \'el\'ement conteneur et n'utilise aucun \verb|document.querySelector| global.
  \item Mesure avec l'onglet \clic{Performance} de Chrome~: d\'eclenche 500 ajouts, trouve la
        t\^ache la plus co\^uteuse, optimise, compare les chiffres.
\end{enumerate}
\end{exercice}

\begin{ressources}
\begin{itemize}
  \item MDN --- R\'ef\'erence DOM --- \url{https://developer.mozilla.org/fr/docs/Web/API/Document\\\\\_Object\\\\\_Model}
  \item web.dev Learn CSS / Learn Performance (Google, gratuit, excellent) --- \url{https://web.dev/learn/}
  \item Chrome DevTools, formation officielle --- \url{https://developer.chrome.com/docs/devtools/learn}
  \item Check-list d'accessibilit\'e --- \url{https://www.a11yproject.com/checklist/}
  \item \link{Can~I~Use}{https://caniuse.com/} pour chaque API ou s\'electeur CSS.
\end{itemize}
\end{ressources}

\index{DOM}\index{navigateur}\index{accessibilite}
